\section{Experimental Setup and Metrics}
\label{sec:setup}

\textbf{Hardware and software.} Training, evaluation and GPU benchmarks ran as non-interactive
Kaggle script jobs on one NVIDIA Tesla T4 (15\,GB usable, fp16, no bf16), PyTorch~2.11 (CUDA~12.8),
Transformers~5.20 and PEFT~0.21. Each job clones the repository at a recorded commit, runs, and
exits so the GPU is released immediately; a short preflight on the longest examples precedes every
training run so that environment or memory errors cost seconds rather than minutes. The application
and the CPU measurements ran on a laptop with an Intel Core i5-1135G7 (4 cores) and 16\,GB RAM,
PyTorch~2.14 (CPU) in fp32. In total all GPU experiments used \TBD{} of the free 30-hour weekly quota.

\textbf{Decoding settings.} All reported predictions use greedy decoding (argmax), at most 64 new
tokens, batch size 16 for evaluation, and \code{repetition\_penalty}$=1.0$. The latter matters:
Qwen's shipped generation configuration enables sampling and a repetition penalty of 1.1, and the
library applies the penalty even when sampling is disabled. We found this when our own decoder
disagreed with library greedy decoding at the third token; with the penalty disabled the two agree
token for token. The zero-shot baseline and the fine-tuned model use the identical prompt, examples
and decoding settings.

\textbf{Classification metrics.} We report accuracy, macro-F1, per-class precision, recall and F1, and
confusion matrices, overall and separately for inclusion and exclusion criteria. Unparseable outputs
count as a separate \emph{invalid} prediction and are always wrong. Macro-F1 averages over the gold
classes present in the evaluated subset; a class with no gold examples (e.g.\ \emph{excluded} in the
test split) has undefined recall and is reported but not averaged, rather than being silently scored
as 0 or 1. Because the overall label space merges the two criterion types, the shared labels
(\emph{not enough information}, \emph{not applicable}) form single classes in the overall figures.

\textbf{Evidence metrics.} Predictions and gold annotations are sets of original sentence IDs.
\emph{Micro} precision/recall/F1 sum true positives, false positives and false negatives over all
records; records with an empty gold set contribute every predicted ID as a false positive.
\emph{Macro} scores average per-record precision, recall and F1 over records with a \emph{non-empty}
gold set (precision is 0 when nothing is predicted). For the 56--70\% of records whose gold set is
empty we report \emph{empty-gold agreement}: the fraction where the model also cites nothing. Absent
evidence is never interpreted as a label.

\textbf{Statistical comparison.} Ablations are compared with a paired bootstrap over validation
annotations (2,000 resamples): for each resample we compute the metric difference between a variant
and the main run, and report the mean difference, its 95\% percentile interval and the fraction of
resamples in which the variant is better.

\textbf{Timing protocol.} Benchmarks use batch size 1 and the same model, weights and inputs for every
mode. Each configuration is run once for warm-up and then repeated (3 times for decoding, 10 for
attention); we report the median. CUDA is synchronized before every clock read, so asynchronous
kernel launches cannot hide work. \emph{Prefill} time is measured up to the first generated token
(including the first argmax), \emph{decode} time covers the remaining tokens, and tokens/s is output
tokens divided by total time. Peak memory is the CUDA allocator peak after resetting its statistics.
For fixed-length comparisons the end-of-sequence token is ignored so every mode generates the same
number of tokens.
